Utility review boards must select grid investment portfolios under hard annual budgets, yet independent project scores ignore portfolio interactions and generic evolutionary operators do not express changes in review terms. BiLo-NSGA embeds budget-aware forward insertion, atomic delete--insert substitution, dependency-aware bonuses, and feasibility recovery within non-dominated sorting; it logs committed local moves and deterministic repair drops. On a reproducible benchmark with 120 candidates, eight scenarios, and 30 seeds per stochastic method, BiLo-NSGA attains pooled mean hypervolume 0.17190, 1.12\% above NSGA-II. Across five stochastic baselines, it has 37 positive mean differences in 40 comparisons, 36 Holm-significant wins, and no significant loss; 16 gaps against two deterministic scoring rules are descriptive and favor BiLo-NSGA. A 1600-neighbor Pareto local-search control is significantly lower in all eight scenarios. Component evidence is asymmetric: removing forward insertion is harmful in three scenarios under the primary family, although the ablation has the higher pooled mean, whereas atomic substitution and standalone deletion are unresolved. Public reliability and MTEP16 backtests provide descriptive external consistency. The evidence supports budget-aware project-level local moves with scenario-dependent forward-side effects; the event log is not a recommendation lineage, and atomic substitution has no demonstrated accuracy gain.
